\section*{Chapter \ref{ch:iv1:expectation-and-indicators} --- Expectation and Indicators}
\markboth{\protect\omTocMark{}{Answers: Expectation and Indicators}}{\protect\omTocMark{}{Answers: Expectation and Indicators}}

\enlargethispage{-4\baselineskip}%
\begin{solution}{iq:iv1:expectation-and-indicators:1}
The expected value is $161/36 \approx 4.47$.

Use the tail-sum formula\index{tail-sum formula}: for a variable taking values $1, 2, \dots$, the mean equals
$\sum_{k \ge 1} P(M \ge k)$, because each outcome $M$ is counted once for every $k$ up to $M$. The higher die is at
least $k$ unless both dice are below $k$, so $P(M \ge k) = 1 - \bigl(\tfrac{k-1}{6}\bigr)^2$. Summing over
$k = 1, \dots, 6$,
\[
E[M] = 6 - \frac{0 + 1 + 4 + 9 + 16 + 25}{36} = 6 - \frac{55}{36} = \frac{161}{36}.
\]
\iqfollow{And with three dice?}{Cubes replace squares: $6 - 225/216 = 119/24 \approx 4.96$.}
\iqfollow{On average, how far apart are the two dice?}{The difference is the higher minus the lower, and
the higher plus the lower is the sum of the dice, which averages 7. So the expected lower die is $7 - 161/36$, and
the difference averages $2 \cdot 161/36 - 7 = 35/18 \approx 1.94$.}
\end{solution}

\begin{solution}{iq:iv1:expectation-and-indicators:2}
The expected total is 21.

The number of rolls $N$ is geometric with mean 6, and each roll has mean $7/2$. Whether you stop after a roll
depends only on the rolls so far, so Wald's identity\index{Wald's identity} applies: $E[\text{total}] = E[N] \cdot
7/2 = 21$. You can check this with first-step analysis\index{first-step analysis}. A roll that is not a six has mean 3, so
\[
S = \tfrac16 \cdot 6 + \tfrac56 \,(3 + S),
\]
which gives $S = 21$.

\iqwrong{``Five earlier rolls averaging 3.5, plus a six, makes 23.5.'' The earlier rolls are known not to be sixes, so
they average 3, not 3.5. Five of them plus the six gives $5 \cdot 3 + 6 = 21$.}
\iqfollow{How many ones do you roll on average before the six?}{Exactly one. Ignore every roll except ones and sixes:
each of those two is equally likely to come next, so the count before the first six is geometric with mean 1.}
\end{solution}

\begin{solution}{iq:iv1:expectation-and-indicators:3}
On average there are exactly 3 such pairs.

Use linearity of expectation with indicators\index{linearity of expectation}\index{indicator variables}. There are 51
neighbouring pairs. Whatever the first card of a pair is, the second is one of the 51 other cards, and 3 of those share its
rank, so each pair matches with probability $3/51$. The total is $51 \cdot 3/51 = 3$. The events overlap and are dependent,
and that makes no difference: linearity needs no independence.

\iqfollow{What is the variance of that count?}{Sum over pairs of pairs. There are 50 overlapping pairs of neighbours, each
matching together with probability $\frac{3}{51}\cdot\frac{2}{50} = \frac{1}{425}$. There are 1\,225 disjoint pairs, each with
probability $\frac{3}{51}\cdot\frac{146}{2450}$, where $\frac{146}{2450} = \frac{2}{50}\cdot\frac{1}{49} +
\frac{48}{50}\cdot\frac{3}{49}$ splits on whether the second pair's first card has the rank just used. Then $\mathrm{Var} = 3 + 2
\cdot \frac{50}{425} + 2\cdot 1225\cdot\frac{3}{51}\cdot\frac{146}{2450} - 9 = \frac{48}{17} \approx 2.82$.}
\iqfollow{How likely is a shuffle with no such pair?}{About 4.5\%, by simulation. A Poisson with mean 3 predicts
$e^{-3} \approx 5\%$, which is close because the many indicators are rare and nearly independent.}
\end{solution}

\enlargethispage{-3\baselineskip}%
\begin{solution}{iq:iv1:expectation-and-indicators:4}
The mean is 1 and the variance is 1, for a deck of any size above one card.

Use indicators\index{indicator variables}: position $i$ matches with probability $1/52$, so the mean is $52 \cdot 1/52 = 1$.
For the variance, write $X^2$ as the sum of the indicators plus the sum over ordered pairs of positions of $I_iI_j$. Positions
$i$ and $j$ both match with probability $\frac{1}{52 \cdot 51}$, and there are $52 \cdot 51$ ordered pairs, so $E[X^2] = 1 + 1 =
2$ and $\mathrm{Var}(X) = 2 - 1 = 1$.

\iqwrong{Treating the positions as independent gives a variance of $52 \cdot \frac{1}{52}\cdot\frac{51}{52} = 51/52$. The small
positive covariance between positions supplies the rest.}
\iqfollow{What is the chance of no match at all?}{Very close to $1/e \approx 0.368$, the fixed-point count being nearly Poisson
with mean 1. This is the chance of a derangement\index{derangements}.}
\end{solution}

\begin{solution}{iq:iv1:expectation-and-indicators:5}
The expected number is $14.7$ rolls.

This is the coupon collector\index{coupon collector}. Split the wait into stages. When you have seen $k$ faces, each roll shows a
new one with probability $(6-k)/6$, so that stage lasts $6/(6-k)$ rolls on average. Adding the stages,
\[
6\Bigl(\tfrac16 + \tfrac15 + \tfrac14 + \tfrac13 + \tfrac12 + 1\Bigr) = 14.7.
\]
\iqfollow{What is the variance?}{The stages are independent geometrics, each with variance $(1-p)/p^2$. Their sum is about
$38.99$, a standard deviation of about $6.2$ rolls. The tail is long because the last face takes 6 rolls on average by itself.}
\iqfollow{What is the expected sum of the rolls?}{Wald's identity\index{Wald's identity} gives $14.7 \cdot 3.5 = 51.45$.}
\iqwhy{The cost is $n \ln n$, not $n$. To see each of 1\,000 instruments at least once in uniformly random samples of the data,
you need about 7\,500 draws, not 1\,000. This is how to size a sample that must touch every bucket.}
\end{solution}

\begin{solution}{iq:iv1:expectation-and-indicators:6}
The mean is 11 and the variance is 5.

Runs are easier to count through changes. The number of runs is one more than the number of neighbouring tosses that
differ. Each of the 20 neighbouring pairs differs with probability $1/2$, so by indicators\index{indicator variables} the mean is
$1 + 20/2 = 11$. The changes are independent, not only pairwise: given any toss, the next one differs with probability $1/2$
whatever came before. So the number of changes is Binomial$(20, 1/2)$, with variance $20/4 = 5$.

\iqfollow{In 100 tosses, how many runs of exactly three heads do you expect?}{A run starting in the middle needs the pattern
THHHT, with probability $1/32$, and there are 96 such starting places. A run at either end needs only HHHT or THHH, with
probability $1/16$ each. The total is $96/32 + 2/16 = 3.125$.}
\iqwhy{This is the runs test\index{runs test} for streaks in returns. For a year of 252 daily signs, the number of runs has mean
$126.5$ and standard deviation about $7.9$, so a count below about 111 is a two-standard-deviation sign of streaks.}
\end{solution}

\begin{solution}{iq:iv1:expectation-and-indicators:7}
The second ace appears at position $106/5 = 21.2$ on average.

Use a symmetry argument\index{symmetry arguments}. The 4 aces split the 48 other cards into 5 gaps, as in \cref{fig:iv1:ace-gaps}.
Each of those 48 cards is equally likely to fall into any gap, because it can sit in any of the 5 places relative to the 4 aces with equal
chance. So each gap holds $48/5 = 9.6$ cards on average. The second ace comes after two gaps and two aces:
\[
2 \cdot \frac{48}{5} + 2 = \frac{106}{5}.
\]
\begin{omfigure}
\begin{tikzpicture}[x=0.9cm,y=0.9cm]
\foreach \a/\b/\l in {0/2/G_1,2.4/4.4/G_2,4.8/6.8/G_3,7.2/9.2/G_4,9.6/11.6/G_5}{
  \fill[omDef!15] (\a,0) rectangle (\b,0.5);
  \draw[omDef] (\a,0) rectangle (\b,0.5);
  \node at ({(\a+\b)/2},0.25) {\small $\l$};
}
\foreach \x/\n in {2.2/1,4.6/2,7/3,9.4/4}{
  \fill[omThm] (\x-0.17,0) rectangle (\x+0.17,0.5);
  \node[above,font=\footnotesize] at (\x,0.5) {ace \n};
}
\draw[decorate,decoration={brace,mirror,amplitude=5pt}] (0,-0.08) -- (4.75,-0.08)
  node[midway,below=5pt,font=\footnotesize] {two gaps and two aces: $2\cdot 9.6 + 2 = 21.2$};
\end{tikzpicture}

\omcaption{The four aces cut the deck into five gaps of equal expected size. The second ace sits after two gaps and the two aces
themselves.}\label{fig:iv1:ace-gaps}
\end{omfigure}

\iqwrong{$2 \cdot 52/5 = 20.8$ forgets that the aces occupy positions too.}
\iqfollow{Where is the last ace on average?}{The same count gives the $k$-th ace at position $k \cdot 53/5$, so the fourth
is at $42.4$.}
\end{solution}

\begin{solution}{iq:iv1:expectation-and-indicators:8}
The expected value is exactly $25.5025$.

Use the tail-sum formula\index{tail-sum formula}. The lowest of the three is at least $k$ when all three dice are at least $k$, which has
probability $\bigl(\tfrac{101-k}{100}\bigr)^3$. Summing from $k = 1$ to $100$ and writing $j = 101 - k$,
\[
E[\min] = \frac{1}{100^3} \sum_{j=1}^{100} j^3 = \frac{1}{100^3}\Bigl(\frac{100 \cdot 101}{2}\Bigr)^2 = \frac{5050^2}{10^6} = 25.5025,
\]
using the fact that the sum of the first $n$ cubes is the square of the sum of the first $n$ integers.

\iqwrong{Using the continuous rule, the lowest of three uniforms sits at a quarter of the range, gives 25. That is close, and it
misses the extra $0.5025$ from the discrete faces. An interviewer asking for the exact value wants the sum.}
\iqfollow{What is the expected highest of the three?}{By the symmetry that sends a face $f$ to $101 - f$, it is $101 - 25.5025 = 75.4975$.}
\end{solution}

\begin{solution}{iq:iv1:expectation-and-indicators:9}
For the die, the expected number of rolls is $1223/324 \approx 3.77$. For the hash table, about 41 keys.

Use the tail-sum formula\index{tail-sum formula}, $E[T] = \sum_{k \ge 0} P(T > k)$. The wait exceeds $k$ when the first $k$ rolls are all different, which has
probability $\frac{6}{6}\cdot\frac{5}{6}\cdots\frac{6-k+1}{6}$. The terms for $k = 0, 1, \dots, 6$ are
\[
1,\quad 1,\quad \tfrac56,\quad \tfrac59,\quad \tfrac{5}{18},\quad \tfrac{5}{54},\quad \tfrac{5}{324},
\]
and their sum is $1223/324$. Note that the wait is at least 2 and at most 7.

For $n$ buckets the same sum is the birthday problem\index{birthday problem}. Its large-$n$ value is about $\sqrt{\pi n/2} + 2/3$, which for
$n = 1024$ gives $40.8$; the exact sum is also $40.8$. That count includes the colliding key.

\iqwhy{Collisions arrive at the square root of the table size, not at the table size. A 32-bit order identifier chosen at random has
its first expected collision after about 82\,000 orders, not 4 billion.}
\end{solution}

\enlargethispage{-3\baselineskip}%
\begin{solution}{iq:iv1:expectation-and-indicators:10}
On average you draw $e \approx 2.718$ numbers.

Let $N$ be the number of draws. Then $N > n$ exactly when the first $n$ numbers sum to at most 1, and that has probability $1/n!$. To see
it, change variables to the running totals $s_1 < s_2 < \dots < s_n$. The map has determinant one, so the probability is the volume of the
region $0 < s_1 < \dots < s_n < 1$, and that region is one of the $n!$ equal orderings of $n$ points in the unit cube. The tail-sum
formula\index{tail-sum formula} gives
\[
E[N] = \sum_{n \ge 0} P(N > n) = \sum_{n \ge 0} \frac{1}{n!} = e.
\]
\iqwrong{``Each number averages one half, so two numbers.'' That ignores the overshoot: the total when you stop is above 1, not equal to 1.}
\iqfollow{What is the average total when you stop?}{By Wald's identity\index{Wald's identity}, $E[N] \cdot \tfrac12 = e/2 \approx 1.359$.}
\iqfollow{What is the variance of $N$?}{$E[N^2] = \sum_{n \ge 0} (2n+1)P(N>n) = 2e + e = 3e$, so the variance is $3e - e^2 \approx 0.766$.}
\end{solution}

\begin{solution}{iq:iv1:expectation-and-indicators:11}
The expected number of further draws is infinite, and yet one further draw is enough half the time.

Let $T$ be the number of further draws. $T > k$ means the first number is the largest of the first $k+1$, and by symmetry
that has probability $1/(k+1)$. So $P(T = 1) = 1/2$, and $P(T \le 3) = 3/4$. The tail-sum formula\index{tail-sum formula}
gives the harmonic series,
\[
E[T] = \sum_{k \ge 0} \frac{1}{k+1} = \infty.
\]
The rare first number that is very large makes for very long waits, and those waits are long enough to outweigh their rarity.

\iqfollow{What if you give up after at most 1\,000 draws?}{The mean is then the harmonic sum $H_{1000} \approx 7.49$. Change the
cap and the answer changes with it, which is the sign of a divergent mean.}
\iqwhy{The average wait from a sample never settles: quadruple the sample and it keeps growing. If you measure ``time to recover the previous high'' from a
backtest with no drift, the sample mean tells you only how long the backtest was. Quote the median or a quantile instead.}
\end{solution}

\begin{solution}{iq:iv1:expectation-and-indicators:12}
The mean is the harmonic number $H_{100} \approx 5.19$, and the variance is about $3.55$.

Use indicators\index{indicator variables}. The $k$-th number is a record when it is the largest of the first $k$, with probability $1/k$, so
$E[R] = \sum_{k=1}^{100} 1/k \approx 5.187$. The large-$n$ form is $\ln n + 0.577$, here about $5.18$.

For the variance you need a stronger fact: the indicators are independent. Knowing how the first $k-1$ numbers are ordered among
themselves says nothing about where the $k$-th ranks among the first $k$, and the record events depend only on those relative ranks. So the
variance is the sum of the individual Bernoulli variances:
\[
\mathrm{Var}(R) = \sum_{k=1}^{100} \frac1k\Bigl(1 - \frac1k\Bigr) \approx 5.187 - 1.635 = 3.55.
\]
\iqwhy{Records in pure noise are common early and rare later: a year of 252 independent daily results has about $6.1$ days that set a new high.
So a strategy that sets six new highs in a year has shown you nothing yet.}
\end{solution}

\enlargethispage{-3\baselineskip}%
\begin{solution}{iq:iv1:expectation-and-indicators:13}
The covariance is $-5/3$ and the correlation is $-1/5$.

Write $A = \sum_i I_i$ and $B = \sum_i J_i$, where $I_i$ marks roll $i$ showing a one, and $J_i$ roll $i$ showing a two. Different rolls are
independent, so only the same roll contributes to the covariance. A single roll cannot be both, so $I_iJ_i = 0$, and
\[
\mathrm{Cov}(I_i, J_i) = 0 - \tfrac16 \cdot \tfrac16 = -\tfrac{1}{36}.
\]
Over 60 rolls this gives $-60/36 = -5/3$. Each count has variance $60 \cdot \tfrac16 \cdot \tfrac56 = 25/3$, so the correlation is
$(-5/3)/(25/3) = -1/5$, whatever the number of rolls.

\iqfollow{Check it another way.}{$A + B$ counts rolls showing a one or a two, which is Binomial$(60, 1/3)$ with variance $40/3$. And indeed
$25/3 + 25/3 - 2\cdot 5/3 = 40/3$.}
\iqwhy{Counts in buckets of a fixed total are negatively correlated by construction. If you split orders across 6 equal venues,
venue shares have correlation $-1/5$ with no economic link between the venues at all, so do not read a negative correlation of that size as
competition.}
\end{solution}

\begin{solution}{iq:iv1:expectation-and-indicators:14}
The expected number of rolls is $3/2$. It is not 3.

Consider stopping at the first roll in $\{1, 3, 5, 6\}$ instead. That stopping time $T$ is geometric with success probability $4/6$, so $E[T] = 3/2$.
Your event ``every roll was even, ending in a six'' is exactly the event that this first stopping roll is a six. The value of the stopping roll is
independent of when it occurs, since each of the four faces is equally likely at any time. So conditioning on it being a six does not change the
distribution of $T$\index{conditioning on the stopping value}.

A direct calculation agrees. The chance of stopping at roll $k$ with every roll even is $(1/3)^{k-1}(1/6)$. Summing over $k$ gives $1/4$, the chance of the condition, and summing $k$
times that term gives $3/8$, so the conditional mean is $(3/8)/(1/4) = 3/2$.

\iqwrong{``Given only evens, the die has three faces, so the wait for a six is geometric with $p = 1/3$, mean 3.'' That is rolling a three-sided die, which
is a different experiment. In ours, a long sequence must also avoid every odd number, so the condition makes long sequences rare.}
\iqfollow{Given only that no one was rolled?}{Stop on $\{1, 6\}$, with success probability $1/3$, so the mean is 3.}
\iqwhy{Conditioning on how a path ended changes how long it seems to have lasted. Take a fair random walk with a stop 10 below and a target 5 above: all trades last 50 steps on average, but the winners last $41.7$ and the losers $66.7$. A backtest report of ``time to target'' from winners alone understates the holding time.}
\end{solution}

\begin{solution}{iq:iv1:expectation-and-indicators:15}
The expected number of loops is $\sum_{k=1}^{100} \frac{1}{2k-1} \approx 3.28$.

Use conditioning on the first step\index{first-step analysis}. Take $n$ noodles, so $2n$ loose ends. Pick one end; its partner is equally likely to be any of the
other $2n - 1$ ends. Exactly one of those is the other end of the same noodle, and then a loop closes. In every other case two noodles merge into
one longer noodle, which behaves exactly like a single noodle. Either way, $n - 1$ noodles remain, so
\[
L_n = L_{n-1} + \text{Bernoulli}\Bigl(\frac{1}{2n-1}\Bigr), \qquad E[L_n] = \sum_{k=1}^{n} \frac{1}{2k-1}.
\]
For large $n$ this is about $\tfrac12 \ln n + 0.98$, and for $n = 100$ that gives $3.28$. Picking two random ends is the same as picking one end
and then a random partner, so nothing is lost by the order we used.

\iqwrong{Guessing a number proportional to $n$. Most ties merge noodles rather than close them, and the loop count grows with the logarithm of $n$.}
\iqfollow{What is the chance that everything ends as one giant loop?}{No tie closes a loop until the very last: $\prod_{k=2}^{100} \frac{2k-2}{2k-1} \approx 0.089$.}
\end{solution}

\begin{solution}{iq:iv1:expectation-and-indicators:16}
The last ant falls after $100/101 \approx 0.990$ minutes on average, and there are $1237.5$ collisions on average.

Ants are identical, so a collision where both turn around looks exactly like two ants passing through each other, as in
\cref{fig:iv1:ants}. Follow these ``ghosts''\index{relabelling identical objects}. Each ghost walks straight to the end it faces and leaves. Its walking time is the distance to that end. Given
a uniform position and a fair direction, that distance is uniform on $[0, 1]$, independently for each ant. The last fall is the maximum of 100
independent uniforms, whose mean is $n/(n+1) = 100/101$.

Each real collision is one crossing of two ghosts. Two ghosts cross exactly when the left one faces right and the right one faces left, with probability
$1/4$, and they then always meet before either falls. So the expected number is $\binom{100}{2}/4 = 1237.5$.

\begin{omfigure}
\begin{tikzpicture}[x=7cm,y=4.8cm]
\draw[->] (0,0) -- (1.08,0) node[right,font=\footnotesize] {position};
\draw[->] (0,0) -- (0,0.92) node[above,font=\footnotesize] {time};
\draw[thick] (1,0) -- (1,0.85);
\node[below,font=\footnotesize] at (0,0) {0};
\node[below,font=\footnotesize] at (1,0) {1};
\draw[omDef,very thick] (0.25,0) -- (0.475,0.225) -- (0,0.7);
\draw[omThm,very thick] (0.7,0) -- (0.475,0.225) -- (1,0.75);
\node[right,font=\footnotesize,xshift=6pt] at (0.475,0.225) {collision};
\node[left,font=\footnotesize,omDef] at (0,0.7) {ant 1 falls at $0.70$};
\node[right,font=\footnotesize,omThm] at (1,0.75) {ant 2 falls at $0.75$};
\node[below,font=\footnotesize,omDef] at (0.25,0) {ant 1};
\node[below,font=\footnotesize,omThm] at (0.7,0) {ant 2};
\end{tikzpicture}

\omcaption{Two ants on the stick, with time running upwards. Each ant turns at the collision and keeps to its own side. Forget the colours and the same picture
shows two straight lines passing through each other, the ghosts. Positions and fall times are the same in both readings.}\label{fig:iv1:ants}
\end{omfigure}

\iqfollow{Does a ghost's fall time belong to the ant it started as?}{No, and the question did not need it. The fall times are correct as a set, not ant by ant:
real ants keep their left-to-right order, so those leaving on the left go in the order they started, and those leaving on the right in reverse order.}
\end{solution}

\begin{solution}{iq:iv1:expectation-and-indicators:17}
On average you turn over $689/45 \approx 15.31$ cards.

Let $A$ be the position of the first ace and $K$ that of the first king. You stop at $\max(A, K)$. Use the max-min identity\index{inclusion-exclusion for maxima}:
\[
\max(A, K) = A + K - \min(A, K).
\]
The first of $m$ special cards in a deck of 52 sits on average at $53/(m+1)$, by the gap symmetry: $m$ special cards split the others into $m+1$ gaps of equal
expected size, and the first special card follows one gap. With 4 aces this is $E[A] = E[K] = 53/5$. With the 8 aces and kings together it is
$E[\min] = 53/9$. So
\[
E[\max(A, K)] = \frac{2 \cdot 53}{5} - \frac{53}{9} = \frac{689}{45}.
\]
\iqwrong{Adding the two waits, $53/5 + 53/5$. They run at the same time, and they overlap by exactly the time to the first special card.}
\iqfollow{A biased three-sided die shows its faces with probabilities $\tfrac12, \tfrac13, \tfrac16$. How many rolls until you have seen all three?}{Inclusion-exclusion
over subsets, with a wait of $1/p$ for any set of total probability $p$: $2 + 3 + 6 - \frac65 - \frac32 - 2 + 1 = 7.3$.}
\iqwhy{This is how to work out the time for both legs of a paired order to fill. If each leg takes 10 seconds on average, independently and memorylessly, then both are done after
$10 + 10 - 5 = 15$ seconds, not 10 and not 20.}
\end{solution}

\begin{solution}{iq:iv1:expectation-and-indicators:18}
The expected number of socks is $256/63 \approx 4.06$.

Use the tail-sum formula\index{tail-sum formula}. Having no pair after $k$ socks means they come from $k$ different pairs. There are $\binom{5}{k}$ ways to choose the pairs and $2^k$ ways
to choose one sock from each, out of $\binom{10}{k}$ draws:
\[
P(T > k) = \frac{2^k \binom{5}{k}}{\binom{10}{k}}, \qquad k = 0, 1, \dots, 5.
\]
These are $1, 1, \tfrac89, \tfrac23, \tfrac{8}{21}, \tfrac{8}{63}$, and their sum is $256/63$. The answer lies between 2 and 6, as it must.

For $n$ pairs the sum has the closed form $4^n / \binom{2n}{n}$, which grows like $\sqrt{\pi n}$. For $n = 5$ the approximation gives $3.96$,
already close.

\iqwrong{Using the birthday formula with $5$ equally likely types of sock gives about $3.51$, which is too low. After the first sock, only 1 of the 9 remaining socks matches it, a chance of $1/9$ and not $1/5$.}
\iqfollow{And with 10 pairs?}{$4^{10}/\binom{20}{10} \approx 5.68$. Doubling the socks multiplies the wait by about $\sqrt{2}$.}
\end{solution}

\begin{solution}{iq:iv1:expectation-and-indicators:19}
On average $62/35 \approx 1.77$ balls are left.

Use an indicator\index{indicator variables} for each ball, and ask when it is still in the urn at the end. A red ball is left over exactly when every blue ball is drawn before it. Look only
at the relative order of this red ball and the 4 blue ones: all 5 orders of ``who comes last'' are equally likely, so this happens with probability $1/5$. Likewise, a blue ball is left
over when all 6 red balls come before it, with probability $1/7$. Summing over balls,
\[
E[\text{left}] = \frac{6}{5} + \frac{4}{7} = \frac{62}{35}.
\]
In general, with $r$ red and $b$ blue, it is $\frac{r}{b+1} + \frac{b}{r+1}$.

\iqwrong{Modelling the draw with fixed probabilities and a geometric wait. The composition of the urn changes with each draw, and the indicators avoid having to track it.}
\iqfollow{What is the chance that the balls left are red?}{They are red exactly when the very last ball in a full draw-down is red, and by symmetry that has probability $6/10 = 3/5$.}
\iqfollow{How many balls do you draw on average?}{$10 - 62/35 = 288/35 \approx 8.23$.}
\end{solution}

\begin{solution}{iq:iv1:expectation-and-indicators:20}
The mean is about $36.60$ and the variance about $9.74$, so the standard deviation is about $3.1$.

Use indicators\index{variance of a count}. Bin $i$ is empty with probability $q_1 = (0.99)^{100}$, so the mean is $100\,q_1 \approx 36.60$, about $100/e$. Bins $i$ and $j$ are both empty with
probability $q_2 = (0.98)^{100}$, since every ball must avoid both. Expanding the square of the sum,
\[
\mathrm{Var}(X) = 100\,q_1 + 100 \cdot 99\, q_2 - (100\,q_1)^2 \approx 9.74.
\]
For large $n$ this tends to $n(e^{-1} - 2e^{-2})$, here $9.72$.

\iqwrong{Treating the count as binomial, $100\,q_1(1 - q_1) \approx 23.2$. That more than doubles the true variance. The bins are negatively correlated, with
$\mathrm{Cov} = q_2 - q_1^2 \approx -0.00136$. If bin 1 is empty, its balls went elsewhere, so other bins are more likely to be filled. This is small per pair, but there are 9\,900 ordered pairs.}
\iqwhy{This gives a quick check on a hash function. With 100 keys and 100 buckets, 50 empty buckets is more than 4 standard deviations above the mean. The binomial variance would have
called it under 3, and you might have accepted the hash.}
\end{solution}

\begin{solution}{iq:iv1:expectation-and-indicators:21}
For 10 numbers, the expected count is about $24.44$ comparisons. In general it is $2(n+1)H_n - 4n$, where $H_n = \sum_{k \le n} 1/k$.

Use one indicator per pair\index{indicator variables}. Name the numbers by their sorted rank, $1, \dots, n$, and fix $i < j$. Two numbers are compared at most once, because a comparison makes one of them the pivot, and a pivot
leaves the list. They stay together until the first pivot chosen from the ranks $i, i+1, \dots, j$. If that pivot is $i$ or $j$, they are compared. If it lies between them, they are separated and never
compared. Every number in that block is equally likely to be the first, so
\[
P(i \text{ and } j \text{ compared}) = \frac{2}{j - i + 1}.
\]
Summing over all pairs gives the formula. For $n = 10$, $2 \cdot 11 \cdot H_{10} - 40 \approx 24.44$.

\iqfollow{Compare with the worst case and with large $n$.}{The worst case is $n(n-1)/2 = 45$ for $n = 10$, when each pivot is the smallest remaining. For large $n$, the mean is about $2n\ln n$, or $1.39\, n \log_2 n$. For $n = 1000$ the exact mean is about 10\,986 comparisons, against 499\,500 in the worst case.}
\end{solution}

\begin{solution}{iq:iv1:expectation-and-indicators:22}
The expected number of deals is exactly 10. For any $n \ge 2$ it is exactly $n$.

The key fact is that a deal of $m$ cards has exactly one match on average, whatever $m$ is, by indicators\index{indicator variables}. So one card leaves per deal on average, and $n$ deals is the natural guess. That is not yet a proof, since the number of deals is random, so prove it by
induction. Let $p_k$ be the chance of $k$ matches in a deal of $n$ cards, with $\sum_k k\,p_k = 1$. Assume $E[R_m] = m$ for all smaller $m$; the case $m = 1$ never arises, since exactly $n - 1$ matches is impossible. By first-step analysis\index{first-step analysis},
\[
E[R_n] = 1 + p_0 E[R_n] + \sum_{k \ge 1} p_k (n - k) = 1 + p_0 E[R_n] + n(1 - p_0) - 1.
\]
So $(1 - p_0)E[R_n] = n(1 - p_0)$, and $E[R_n] = n$.

\iqwrong{Thinking this needs the derangement probabilities. It does not: they cancel.}
\iqfollow{How many individual cards are dealt in total?}{Each deal of $m$ cards costs $m$. Using the same first step with $E[k] = 1$ and $E[k^2] = 2$, the total $S_n = n^2/2 + n$ satisfies the recursion. For $n = 10$ that is 60 cards.}
\end{solution}

\begin{solution}{iq:iv1:expectation-and-indicators:23}
The expected value of the final number is $e/2 - 1 \approx 0.359$.

Compute by summing over how long the increasing sequence lasts. Fix $n \ge 1$ and a value $y$. The sequence stops at draw $n+1$, with that draw equal to $y$, when the first $n$ are increasing and the $n$-th exceeds $y$. The first $n$ are increasing with probability $1/n!$. They are
increasing and all below $y$ with probability $y^n/n!$. So the probability we want is $(1 - y^n)/n!$, and
\[
E[\text{final}] = \sum_{n \ge 1} \int_0^1 y\,\frac{1 - y^n}{n!}\,dy = \sum_{n \ge 1} \frac{1}{n!}\Bigl(\frac12 - \frac{1}{n+2}\Bigr).
\]
The first part is $(e-1)/2$. For the second, $\sum_{n \ge 0} \frac{1}{n!(n+2)} = \int_0^1 x e^x\,dx = 1$. Dropping the $n = 0$ term, which is $\tfrac12$, leaves $\tfrac12$. Hence
\[
E[\text{final}] = \frac{e - 1}{2} - \frac12 = \frac{e}{2} - 1.
\]
\iqfollow{And the number just before the final one, the highest in the sequence?}{The same method gives $\sum_{n \ge 1} \frac{1}{(n-1)!(n+2)} = \int_0^1 x^2 e^x\,dx = e - 2 \approx 0.718$.}
\iqfollow{How can you check the answer?}{The number of draws has mean $e$, as for any increasing run, and it is a stopping time. So Wald's identity\index{Wald's identity} says the total averages $e/2$. A simulation confirms the total, the $0.359$ and the $0.718$.}
\end{solution}

\begin{solution}{iq:iv1:expectation-and-indicators:24}
The expected number is $1 + \tfrac12 + \tfrac13 + \tfrac14 + \tfrac15 + \tfrac16 = 49/20 = 2.45$. In general, with $n$ faces, it is the harmonic number $H_n$.

Label the faces by the order in which they first appear, and take the $k$-th new face, $f$. At the moment $f$ first appears, $6 - k$ faces are still unseen. Then $f$ ends up appearing exactly once if it does not come up again before all the unseen faces have appeared.

Ignore every roll that shows neither $f$ nor one of those unseen faces\index{symmetry arguments}. The remaining rolls are uniform on $7 - k$ faces and independent of the past. By symmetry each of these $7 - k$ faces is equally likely to be the last of them to appear in the future rolls, so $f$ is last with probability $1/(7-k)$. If $f$ is last, every unseen face comes before its next appearance. So the $k$-th new face is a singleton with probability $1/(7-k)$, and by linearity\index{linearity of expectation}
\[
E[\text{singletons}] = \sum_{k=1}^{6} \frac{1}{7-k} = H_6 = 2.45.
\]
\iqwrong{Using a Poisson count at the expected stopping time: with $14.7$ rolls, each face appears once with chance about $(14.7/6)\,e^{-14.7/6}$, so the total is about $1.27$. The stopping time depends on the counts: the rule ends the game as soon as the last new face appears, so that face is always a singleton. Treating the time as fixed misses this.}
\iqfollow{With a 52-card deck, sampling with replacement until you have seen all cards?}{$H_{52} \approx 4.54$. It grows very slowly, while the total wait grows like $n \ln n$.}
\end{solution}

\begin{solution}{iq:iv1:expectation-and-indicators:25}
The expected number of rolls is about $24.13$.

Poissonisation\index{Poissonisation} makes the faces independent. Imagine rolls arriving as a Poisson process of rate one per unit time. Each face then appears as an independent Poisson process of rate $1/6$. The expected stopping time in clock time equals the expected number of rolls, because each roll takes one unit on average and the stopping rule looks only at the rolls.

By time $t$, with $u = t/6$, a given face has appeared at least twice with probability $1 - e^{-u}(1 + u)$. All six have done so with the sixth power of this. Integrating the tail,
\[
E[\text{rolls}] = \int_0^\infty P(\tau > t)\,dt = 6 \int_0^\infty \Bigl[1 - \bigl(1 - e^{-u}(1 + u)\bigr)^6\Bigr] du.
\]
Expand by the binomial theorem. Each term is an elementary integral:
\[
E[\text{rolls}] = 6 \sum_{k=1}^{6} (-1)^{k+1} \binom{6}{k} \int_0^\infty e^{-ku}(1 + u)^k\,du,
\]
where the $k = 1$ integral is $2$, the $k = 2$ integral is $5/4$, and so on. The sum comes to $24.13$.

An alternative that you can run as a table is a Markov chain on $(a, b)$, the numbers of faces seen zero times and once, with $E(a,b) = \bigl[6 + a E(a-1, b+1) + b E(a, b-1)\bigr]/(a+b)$ and $E(0,0)=0$.

\iqwrong{Doubling the single-collection wait to $29.4$. Collecting the first full set leaves duplicates, which count towards the second. Only $9.43$ more rolls are needed on average beyond the $14.7$.}
\end{solution}
